\documentclass[12pt]{article}
\usepackage[utf8]{inputenc}
\usepackage[T1]{fontenc}
\usepackage[margin=1in]{geometry}
\usepackage{setspace}
\usepackage{xcolor}
\usepackage[colorlinks=true, urlcolor=blue!60!black, linkcolor=blue!60!black]{hyperref}
\setlength{\parindent}{2em}
\setlength{\parskip}{1.25em}
\renewcommand{\baselinestretch}{1.0}
\begin{document}
\begin{center}
{\LARGE \textit{United States v. Cruikshank}, 92 U.S. 542 (1876)\par}
\vspace{0.8em}
{\large Case Brief --- Four-Page Edition\par}
\end{center}
\vspace{1em}

\section{Memory Jogger}

Overturning the federal conspiracy convictions that arose from the Colfax Massacre, the Supreme Court held that the First and Second Amendments bind only Congress and that the Fourteenth Amendment bars only state---not private---action, severely curbing federal power to prosecute racial terror during Reconstruction (\href{https://supreme.justia.com/cases/federal/us/92/542/}{Justia}).

\section{Detailed Case Facts}

After Louisiana's disputed 1872 election, Black Republican freedmen guarded the Grant Parish courthouse in Colfax against a feared Democratic takeover. On Easter Sunday, April 13, 1873, an armed white mob attacked with rifles and a cannon and murdered scores of Black men, many after they surrendered; estimates range from roughly 60 to 150 killed. With no state prosecution, federal prosecutors indicted more than one hundred men under Section~6 of the Enforcement Act of 1870, which made it a federal felony to ``band or conspire together'' to hinder any citizen's ``free exercise and enjoyment of any right or privilege granted or secured to him by the Constitution or laws of the United States.'' Sixteen counts charged conspiracies against Levi Nelson and Alexander Tillman's rights to assemble, bear arms, enjoy life and liberty, equal protection, and vote; sixteen more charged murder.

\section{Procedural History}

At a first trial in the Circuit Court for the District of Louisiana, one defendant was acquitted and the jury deadlocked on the rest. On retrial in June 1874 before Judge Woods and Justice Bradley, riding circuit, William Cruikshank, John Hadnot, and William Irwin were convicted on the sixteen conspiracy counts and acquitted of murder. Bradley would have arrested judgment; Woods disagreed, and the division was certified to the Supreme Court, which affirmed the arrest of judgment on March 27, 1876.

\section{Judicial Votes}

Arrest of judgment affirmed, \textbf{9--0} in the judgment, \textbf{8--1} on the reasoning.
\begin{itemize}
\item \textbf{Majority (8):} Chief Justice Waite, joined by Justices Swayne, Miller, Davis, Field, Strong, Bradley, and Hunt.
\item \textbf{Concurring in the judgment (styled a dissent):} Justice Clifford, on pleading grounds only.
\end{itemize}

\section{Holding}

None of the sixteen counts stated an offense under Section~6. Congress may protect only rights of national citizenship. The First and Second Amendments restrict only Congress; only assembly for a federal purpose, such as petitioning Congress, is a national right. The Fourteenth Amendment bars state encroachment only and adds nothing to one private citizen's rights against another. The Fifteenth Amendment confers only freedom from racial discrimination in voting, and counts omitting racial motivation were defective. Counts charging a conspiracy against ``every, each, all, and singular'' rights were fatally vague. Because the government pressed only three counts, the broad constitutional rulings exceeded what was necessary; \textit{De Jonge} (1937) and \textit{McDonald} (2010) later incorporated assembly and arms rights, while the state-action rule survives.

\section{Analysis of Opinions}

\textbf{Majority (Waite).} Building on the \textit{Slaughter-House Cases}, Waite stressed dual citizenship and delegated federal power, then went count by count. Assembly predated the Constitution and, under \textit{Barron v. Baltimore}, the First Amendment restricts only Congress; the indictment alleged no federal purpose. The Second Amendment leaves citizens to state police power for protection. Murder and false imprisonment are state offenses; the Fourteenth Amendment guards only against state denial, and the civil-rights counts did not allege racial motive. Under \textit{Minor v. Happersett} and \textit{Reese}, voting in state elections derives from the States; though the Court ``may suspect that race was the cause of the hostility,'' the counts did not aver it or specify federal elections. The generic counts failed the Sixth Amendment's demand that an indictment ``descend to particulars.''

\textbf{Clifford (concurring in the judgment).} He agreed judgment must be arrested ``but for reasons quite different from those given by the court.'' Federal courts have no common-law criminal jurisdiction, and counts 13, 14, and 16, the only ones argued, lacked particularity about the right invaded or the elections at issue. Unlike Waite, he recognized Congress's power to enforce the Privileges or Immunities Clause and did not reach the constitutional challenges.

\section{Examples: Future Applications}

\begin{enumerate}
\item \textbf{Same-Side: Private Zoning-Dispute Assault.} Private individuals assault a rival over a municipal ordinance, with no state involvement, racial animus, or federal nexus. The remedy lies in state law.
\item \textbf{Same-Side: Generic \S~241 Indictment.} An indictment alleges a conspiracy to deprive a victim of ``all rights'' secured by federal law without naming one. Under the pleading rule both Waite and Clifford endorsed, it must be dismissed.
\item \textbf{Opposite-Side: Petitioning Congress or Federal Elections.} Paramilitaries disperse citizens petitioning Congress or intimidate Black voters in a House election. Even \textit{Cruikshank} treats these as national rights Congress may protect against private conspiracies.
\item \textbf{Opposite-Side: Officials Conspiring with a Mob.} A sheriff releases civil-rights workers to Klansmen who murder them. State action is present, as in \textit{United States v. Price} and \textit{Guest} (1966).
\item \textbf{Fence-Sitter: Systematic Failure to Prosecute.} State prosecutors and juries refuse to punish white mobs terrorizing Black citizens. \textit{Cruikshank} left open whether such systematic omission is itself state action under Section~5.
\end{enumerate}

\section{Critique}

\textbf{Originalist and textualist critiques.} Randy Barnett, Evan Bernick, and Ilya Shapiro, echoing Justice Thomas in \textit{McDonald}, argue that the Reconstruction framers meant the Privileges or Immunities Clause to secure fundamental rights, including arms, against the States, and that citing \textit{Barron} as if nothing had changed defied original meaning. Some textualists defend the state-action rule because Section~1 begins ``No State shall.''

\textbf{Progressive and historical critiques.} James Gray Pope argues that \textit{Cruikshank} ``belongs at the heart of the American constitutional canon'' because it immunized white supremacist terror and paved the way for Jim Crow. Critics fault the formalism of merely ``suspecting'' racial motive in a racial massacre and the neglect of the view that failing to protect Black citizens denies equal protection.

\section{Key Quotations}

\begin{quote}
``This is not a right granted by the Constitution. Neither is it in any manner dependent upon that instrument for its existence. The second amendment declares that it shall not be infringed, but this, as has been seen, means no more than that it shall not be infringed by Congress.''
\end{quote}
\noindent --- Chief Justice Waite, Opinion of the Court.

\begin{quote}
``The Fourteenth Amendment prohibits a State from depriving any person of life, liberty, or property without due process of law, but this adds nothing to the rights of one citizen as against another. It simply furnishes an additional guaranty against any encroachment by the States upon the fundamental rights which belong to every citizen as a member of society.''
\end{quote}
\noindent --- Chief Justice Waite, Opinion of the Court.

\begin{quote}
``The same person may be at the same time a citizen of the United States and a citizen of a State, but his rights of citizenship under one of those governments will be different from those he has under the other.''
\end{quote}
\noindent --- Chief Justice Waite, Opinion of the Court.

\end{document}
